\documentclass[pdflatex,sn-nature]{sn-jnl}% Same class and reference style as main.tex

\usepackage{graphicx}%
\usepackage{amsmath,amssymb,amsfonts}%
\usepackage{xcolor}%
\usepackage{textcomp}%
\usepackage{booktabs}%
\usepackage{array}%

\unnumbered% Nature-style unnumbered section heads (sn-jnl macro)
\renewcommand{\labelitemiii}{$\circ$}
% Handwriting space between outline bullets; set to 0pt to turn off.
\newlength{\writinggap}\setlength{\writinggap}{5\baselineskip}
\makeatletter
\newcommand{\writingspace}{\g@addto@macro\itemargs{\itemsep\writinggap\topsep\writinggap}}
\makeatother
\writingspace

\graphicspath{{fig/supp/}}

\raggedbottom

% Supplementary float labels: "Supplementary Fig. S1", "Supplementary Table S1".
\renewcommand{\figurename}{Supplementary Fig.}
\renewcommand{\thefigure}{S\arabic{figure}}
\renewcommand{\tablename}{Supplementary Table}
\renewcommand{\thetable}{S\arabic{table}}
\providecommand{\theHfigure}{}
\renewcommand{\theHfigure}{S\arabic{figure}}
\providecommand{\theHtable}{}
\renewcommand{\theHtable}{S\arabic{table}}

% Supplementary References continue the main list's numbering.
% Set this to the number of distinct \cite keys in main.tex (37 on 2026-10-09).
% Note from Claude: main-with-methods-notes.tex had 39 distinct keys on 2026-10-23 while Methods was still being written. Recount once Methods is final.
\newcommand{\SImainrefcount}{39}
% natbib's thebibliography zeroes NAT@ctr when it opens, so the offset goes after it.
\makeatletter
\let\SIorigthebibliography\thebibliography
\renewcommand{\thebibliography}[1]{%
  \SIorigthebibliography{#1}%
  \global\c@NAT@ctr=\SImainrefcount\relax}
\makeatother

\begin{document}

\title{Supplementary Information for: \textcolor{red}{[TITLE TBD]}}

\author*[1]{\fnm{Eva Maxfield} \sur{Brown}}\email{evamxb@uw.edu}

\author[1]{\fnm{Nicholas} \sur{Weber}}\email{nmweber@uw.edu}

\affil*[1]{\orgdiv{Information School}, \orgname{University of Washington}, \orgaddress{\city{Seattle}, \state{WA}, \country{USA}}}

\maketitle

%%%%%%%%%%%%%%%%%%%%%%%%%%%%%%%%%%%%%%%%%%%%%%%%%%%%%%%%%%%%%%%%%%%%%%%%%%%%%%%%
% Supplementary Notes
%%%%%%%%%%%%%%%%%%%%%%%%%%%%%%%%%%%%%%%%%%%%%%%%%%%%%%%%%%%%%%%%%%%%%%%%%%%%%%%%

\section*{Supplementary Note 1: Co-authorship network connectivity before and after mining}\label{snote:network}

% Sources: outputs/seed_vs_full_network_comparison.csv and outputs/coauthorship_network_summary.json (not the .prev copy).
Our network mining process grows outward from a known co-authorship network similar to a snowball sampling process.
As a result, the network becomes increasingly connected and concentrated in a single large component.
To evaluate network connectivity, we compare the co-authorship network created from the seed pairs alone to the co-authorship network created from the seed and mined pairs together.
In each network, researchers are nodes and an undirected edge joins two researchers who share an RS-Graph article.
Single-author articles add no edges, so we exclude researchers without any co-authorship edge, which removes 3,834 researchers from the seed-only network and 3,132 from the seed and mined network.
The largest connected component holds 84.9\% of researchers in the seed-only network (354,036 of 417,082) and 93.6\% in the seed and mined network (590,147 of 630,651).
Over the same change, the number of connected components falls from 14,145 to 9,387, and the second-largest component in the seed and mined network has only 49 researchers.
We expect this increase by construction, because we find every mined pair through a researcher who already authored an article in RS-Graph.

\section*{Supplementary Note 2: ARMM input format}\label{snote:armminput}

% Source: sci-soft-models binary_article_repo_em/constants.py (OPTIMIZED_MODEL_STR_INPUT_TEMPLATE, truncation lengths, TRUNC_ELEMENT); main.py _prepare_input and training.py use the same template, and missing values become the string "None" in both.
We pass each candidate article-repository pair to the ARMM as a single XML string, built from the template below by replacing each placeholder in braces with the pair's value.
The same template builds the training, evaluation, and mining inputs.
We truncate abstracts at 3,000 characters and READMEs at 6,000 characters, and we append \texttt{<TRUNC/>} to any text we truncate.
We write a missing value as the string \texttt{None}.

\begingroup
\makeatletter\def\verbatim@font{\normalfont\ttfamily\fontsize{6.8}{8.5}\selectfont}\makeatother
\begin{verbatim}
<document-repository-comparison>
    <document>
        <title>{document_title}</title>
        <doi>{document_doi}</doi>
        <publication_date>{document_publication_date}</publication_date>
        <abstract_content_truncated>{document_abstract_content_truncated}</abstract_content_truncated>
    </document>
    <repository>
        <owner>{repository_owner}</owner>
        <name>{repository_name}</name>
        <description>{repository_description}</description>
        <primary_language>{repository_primary_language}</primary_language>
        <is_fork>{repository_is_fork}</is_fork>
        <creation_datetime>{repository_creation_datetime}</creation_datetime>
        <last_pushed_datetime>{repository_last_pushed_datetime}</last_pushed_datetime>
        <stargazers_count>{repository_stargazers_count}</stargazers_count>
        <forks_count>{repository_forks_count}</forks_count>
        <open_issues_count>{repository_open_issues_count}</open_issues_count>
        <commits_count>{repository_commits_count}</commits_count>
        <size_kb>{repository_size_kb}</size_kb>
        <topics>{repository_topics}</topics>
        <readme_content_truncated>{repository_readme_content_truncated}</readme_content_truncated>
    </repository>
</document-repository-comparison>
\end{verbatim}
\endgroup

\clearpage

%%%%%%%%%%%%%%%%%%%%%%%%%%%%%%%%%%%%%%%%%%%%%%%%%%%%%%%%%%%%%%%%%%%%%%%%%%%%%%%%
% Supplementary Tables
%%%%%%%%%%%%%%%%%%%%%%%%%%%%%%%%%%%%%%%%%%%%%%%%%%%%%%%%%%%%%%%%%%%%%%%%%%%%%%%%

% Table S1 (sideways: six text columns wrap to 4--8 lines at text width)
% Wattanakriengkrai et al. facts checked against the arXiv full text (2004.00199v3): 20,278 READMEs matching DOI, BibTeX, or arXiv patterns, repositories created 2014--2018 in C, C++, Java, JavaScript, Python, PHP, and Ruby; machine learning the most frequent topic; more than half of referenced papers do not link back. Venue confirmed as JSS 183, 111117 (DOI 10.1016/j.jss.2021.111117), matching main.bib.
% OpenAIRE cell is qualitative only (Eva, 2026-10-23). Source: memory 2026-10-14-openaire-stage2-isrelatedto.
\begin{sidewaystable}
\caption{\textbf{Prior resources linking scientific articles, code, and software use.} Resources are listed with the scope their own titles and descriptions report. ``Official pairs'' means links between an article and the repository holding its own code, as defined in the main text.}\label{stab:prior}
\footnotesize
\setlength{\tabcolsep}{3pt}
\renewcommand{\arraystretch}{1.1}
\begin{tabular}{@{}>{\raggedright\arraybackslash}p{84pt}>{\raggedright\arraybackslash}p{124pt}>{\raggedright\arraybackslash}p{92pt}>{\raggedright\arraybackslash}p{132pt}>{\centering\arraybackslash}p{40pt}>{\raggedright\arraybackslash}p{50pt}@{}}
\toprule
Resource & What it links & Fields & Official article-repository pairs & Identities\textsuperscript{a} & Software-use views \\
\midrule
\multicolumn{6}{@{}l}{\textit{Article--code links}}\\
\addlinespace[2pt]
Papers with Code & Papers to code implementations & Mostly CS and ML & Author-provided links mixed with community reimplementations & -- & -- \\ \addlinespace
Linked\-Papers\-With\-Code~\cite{farber2023linked} & Papers with Code as an RDF knowledge graph & ML & As in Papers with Code & -- & -- \\ \addlinespace
Science Data Lake~\cite{wilinski2026science} & Eight scholarly sources joined by DOI, including Papers with Code & All fields (bibliometric), with paper--code links from Papers with Code only & As in Papers with Code & -- & -- \\ \addlinespace
OpenAIRE Graph~\cite{manghi2019openaire} & Publication--software relations within a scholarly knowledge graph & All fields & Not distinguished, because inferred links are all typed \texttt{IsRelatedTo} and do not separate an article's own code from software it used & -- & -- \\ \addlinespace
Wattanakriengkrai et al.~\cite{wattanakriengkrai2022github} & 20,278 GitHub repositories created 2014--2018 whose READMEs reference academic papers & Mixed, with ML the most common topic & Not tested, since a README may cite any paper and more than half of cited papers do not link back & -- & -- \\ \addlinespace
Trisovic et al.~\cite{trisovic2022large} & Replication datasets in Harvard Dataverse and their R code & Studies depositing replication data in Harvard Dataverse & Code deposited with the study & -- & -- \\
\midrule
\multicolumn{6}{@{}l}{\textit{Software mentions}}\\
\addlinespace[2pt]
Softcite dataset~\cite{du2021softcite} & Software mentions in articles & Biomedical and economic research & -- & -- & Mentions \\ \addlinespace
SoMeSci~\cite{schindler2021somesci} & Gold-standard knowledge graph of software mentions & Scientific articles & -- & -- & Mentions \\ \addlinespace
CZ Software Mentions~\cite{istrate2022large} & Software mentions in biomedical papers, some disambiguated and linked to the mentioned software's repository & Biomedical & -- & -- & Mentions \\ \addlinespace
SoftCite-2025 extractions~\cite{howison_2025_15149379} & Software mentions across the open-access literature & All fields (open access) & -- (used here as a seed source) & -- & Mentions \\
\midrule
RS-Graph v1~\cite{brown2026code} & 163,292 article-repository pairs from Papers with Code, PLOS, JOSS, and SoftwareX & Mostly CS, since most pairs come from Papers with Code & Stated links & Yes & -- \\ \addlinespace
\textbf{RS-Graph (this work)} & \textbf{309,959 article-repository pairs from five seed sources plus mining} & \textbf{26 fields} & \textbf{Stated links, classifier-derived seeds, and mined pairs, each with its route and confidence} & \textbf{Yes} & \textbf{Mentions, imports, and dependencies} \\
\botrule
\end{tabular}
\footnotetext{\textsuperscript{a}Researcher-developer-account identities. CS, computer science. ML, machine learning. --, not provided.}
\end{sidewaystable}

% Table S2
\begin{table}[p]
\caption{\textbf{RS-Graph data dictionary.} Tables and key columns in RS-Graph (\texttt{rs\_graph/db/models.py}), grouped by node and edge type. Every table also carries an integer \texttt{id} and \texttt{created\_datetime} / \texttt{updated\_datetime} collection timestamps. Fields marked \textsuperscript{\dag} are omitted from the public release. Seed pairs have an empty \texttt{iteration}. Mined pairs have the \texttt{dataset\_source} \texttt{snowball-sampling-discovery} and their mining round in \texttt{iteration}.}\label{stab:dictionary}
\footnotesize
\setlength{\tabcolsep}{3pt}
\begin{tabular}{@{}>{\raggedright\arraybackslash}p{94pt}>{\raggedright\arraybackslash}p{92pt}>{\raggedright\arraybackslash}p{174pt}@{}}
\toprule
Table & Holds & Key columns \\
\addlinespace[1pt]
\multicolumn{3}{@{}l}{\textbf{Articles}}\\
\texttt{document} & One row per article & \texttt{doi}, \texttt{open\_alex\_id}, \texttt{title}, \texttt{publication\_date}, \texttt{cited\_by\_count}, \texttt{fwci}, \texttt{citation\_normalized\_percentile}, \texttt{document\_type}, \texttt{is\_open\_access}, \texttt{primary\_location\_id}, \texttt{best\_open\_access\_location\_id} \\
\texttt{document\_alternate\_doi} & Alternate DOIs (e.g., a preprint's DOI) & \texttt{document\_id}, \texttt{doi} \\
\texttt{document\_abstract}\textsuperscript{\dag} & Article abstracts & \texttt{document\_id}, \texttt{content} \\
\texttt{location}, \texttt{source} & Where an article is published or hosted & \texttt{landing\_page\_url}, \texttt{pdf\_url}, \texttt{is\_open\_access}, \texttt{license}, \texttt{version}; venue \texttt{name}, \texttt{source\_type}, \texttt{host\_organization\_name} \\
\texttt{topic}, \texttt{document\_topic} & OpenAlex topics and their fields, subfields, and domains & \texttt{name}, \texttt{field\_name}, \texttt{subfield\_name}, \texttt{domain\_name}; per-document \texttt{score} \\
\texttt{funder}, \texttt{funding\_instance}, \texttt{document\_funding\_\allowbreak instance} & Funding acknowledged by each article & funder \texttt{name}, \texttt{award\_id} \\
\addlinespace[1pt]
\multicolumn{3}{@{}l}{\textbf{Researchers}}\\
\texttt{researcher} & One row per OpenAlex author record & \texttt{open\_alex\_id}, \texttt{orcid}, \texttt{name}, \texttt{works\_count}, \texttt{cited\_by\_count}, \texttt{h\_index}, \texttt{i10\_index}, \texttt{two\_year\_mean\_citedness} \\
\texttt{institution}, \texttt{document\_contributor\_\allowbreak institution} & Author affiliations on each article & \texttt{name}, \texttt{country\_code}, \texttt{institution\_type}, \texttt{ror} \\
\addlinespace[1pt]
\multicolumn{3}{@{}l}{\textbf{Repositories}}\\
\texttt{code\_host} & Code host (GitHub) & \texttt{name} \\
\texttt{repository} & One row per repository & \texttt{owner}, \texttt{name}, \texttt{description}, \texttt{is\_fork}, \texttt{forks\_count}, \texttt{stargazers\_count}, \texttt{open\_issues\_count}, \texttt{commits\_count}, \texttt{size\_kb}, \texttt{topics}, \texttt{primary\_language}, \texttt{license}, \texttt{processed\_at\_sha}, \texttt{creation\_datetime}, \texttt{last\_pushed\_datetime} \\
\texttt{repository\_readme} & README text & \texttt{repository\_id}, \texttt{content} \\
\texttt{repository\_language} & Bytes of code per language & \texttt{language}, \texttt{bytes\_of\_code} \\
\texttt{repository\_file} & File tree & \texttt{path}, \texttt{tree\_type}, \texttt{bytes\_of\_code} \\
\addlinespace[1pt]
\multicolumn{3}{@{}l}{\textbf{Developer accounts}}\\
\texttt{developer\_account} & One row per account & \texttt{username}, \texttt{name}\textsuperscript{\dag}, \texttt{email}\textsuperscript{\dag} \\
\addlinespace[1pt]
\multicolumn{3}{@{}l}{\textbf{Edges}}\\
\texttt{document\_contributor} & Authorships & \texttt{researcher\_id}, \texttt{document\_id}, \texttt{position}, \texttt{is\_corresponding} \\
\texttt{repository\_contributor} & Code contributions & \texttt{repository\_id}, \texttt{developer\_account\_id} \\
\texttt{researcher\_developer\_\allowbreak account\_link} & Researcher-developer-account identities & \texttt{researcher\_id}, \texttt{developer\_account\_id}, \texttt{predictive\_model\_name}, \texttt{predictive\_model\_version}, \texttt{predictive\_model\_confidence}, \texttt{last\_snowball\_processed\_datetime} \\
\texttt{document\_repository\_link} & Article-repository pairs & \texttt{document\_id}, \texttt{repository\_id}, \texttt{dataset\_source\_id} (route), \texttt{iteration} (mining round), \texttt{predictive\_model\_name}, \texttt{predictive\_model\_version}, \texttt{predictive\_model\_confidence} \\
\texttt{dataset\_source} & Route names & \texttt{name} \\
\addlinespace[1pt]
\multicolumn{3}{@{}l}{\textbf{Software-use views}}\\
\texttt{document\_software\_\allowbreak mention} & Software mentions from SoftCite-2025 & \texttt{softcite\_mention\_id}, \texttt{software\_name}, \texttt{software\_name\_normalized}, \texttt{mention\_context} \\
\texttt{repository\_import} & Imported third-party libraries & \texttt{software\_name}, \texttt{software\_name\_normalized}, \texttt{file\_paths} \\
\texttt{repository\_dependency} & Declared dependencies & \texttt{software\_name}, \texttt{software\_name\_normalized}, \texttt{version\_spec}, \texttt{ecosystem}, \texttt{manifest\_paths}, \texttt{dependency\_type} \\
\botrule
\end{tabular}
% Source for the seed/mined rule: publication scripts treat iteration IS NULL as seed (figure2_dataset_coverage.py, figure3_development_characteristics.py); utils.SOURCE_DISPLAY_NAMES maps "snowball-sampling-discovery" to Mined.
\end{table}

% Table S3
\begin{sidewaystable}
\caption{\textbf{Characteristics of seed and mined pairs.} Each column uses the same population as the corresponding statistic in Results: FWCI over all articles; license, manifest, and Python share over repositories linked to articles published 2008--2025; timing per pair. Python share here is pooled over 2008--2025 (53.6\% for all pairs), while Results pools 2012--2025 (53.7\%). No repository is linked through both a seed pair and a mined pair.}\label{stab:seedmined}
\setlength{\tabcolsep}{3pt}
\begin{tabular}{@{}>{\raggedright\arraybackslash}p{88pt}*{9}{>{\raggedleft\arraybackslash}p{39pt}}>{\raggedleft\arraybackslash}p{52pt}@{}}
\toprule
Group & Pairs & Articles with FWCI & Median FWCI & Articles with FWCI $>$ 1 & Permissive license & No license & Manifest (2008--2025) & Manifest (2025) & Python (known language) & Median days created before / last push after publication \\
\midrule
All pairs & 309,959 & 57.6\% & 2.04 & 65.1\% & 35.8\% & 49.9\% & 34.0\% & 42.5\% & 53.6\% & 96 / 117 \\
Seed (all sources) & 196,107 & 54.2\% & 2.47 & 68.7\% & 39.4\% & 46.6\% & 38.6\% & 46.4\% & 61.0\% & 97 / 164 \\
Seed: Papers with Code & 168,456 & 51.4\% & 2.70 & 70.4\% & 41.1\% & 46.6\% & 40.8\% & 47.8\% & 65.4\% & 81 / 163 \\
Seed: other four sources & 27,651 & 70.6\% & 1.63 & 61.5\% & 29.1\% & 46.1\% & 24.6\% & 28.2\% & 31.6\% & 268 / 176 \\
Mined (rounds 1--5) & 113,852 & 63.3\% & 1.59 & 59.9\% & 29.5\% & 55.5\% & 26.1\% & 38.5\% & 40.4\% & 95 / 51 \\
Mined, excluding Computer Science & 64,727 & 68.6\% & 1.76 & 62.2\% & 27.6\% & 55.5\% & 21.9\% & 30.2\% & 30.4\% & 118 / 49 \\
Seed, excluding Computer Science & 82,607 & 61.3\% & 2.38 & 68.4\% & 35.2\% & 47.2\% & 31.6\% & 39.4\% & 46.7\% & 166 / 142 \\
\botrule
\end{tabular}
\footnotetext{Median FWCI includes articles with zero citations. Percentages of articles with FWCI are of all articles in the group; the median and the share above 1 are among articles with an FWCI.}
\end{sidewaystable}

% Table S4
\begin{table}[p]
\caption{\textbf{Repository timing relative to article publication, by pair source.} Median days from repository creation to article publication, and from article publication to the repository's last push, over all pairs in the seeded or high-precision subset. Negative values mean the event came after (creation) or before (last push) publication.}\label{stab:timing}
\begin{tabular}{@{}>{\raggedright\arraybackslash}p{70pt}>{\raggedleft\arraybackslash}p{40pt}*{3}{>{\raggedleft\arraybackslash}p{70pt}}@{}}
\toprule
Source & Pairs & Median days created before publication & Median days of last push after publication & Pairs with no push after publication \\
\midrule
JOSS & 2,969 & 561 & 718 & 8.3\% \\
SoftCite-2025 & 16,653 & 236 & 295 & 27.8\% \\
PLOS & 7,361 & 246 & $-$34 & 64.8\% \\
SoftwareX & 668 & 216.5 & 144 & 28.7\% \\
Papers with Code & 168,456 & 81 & 163 & 22.4\% \\
Mined & 113,852 & 95 & 51 & 32.5\% \\
All pairs & 309,959 & 96 & 117 & 27.3\% \\
\botrule
\end{tabular}
\end{table}

% Table S5 (sideways: five name-and-count columns do not fit the text width)
% Sources: upper block outputs/top_software_by_view.csv (mentions_as_is, mentions_stoplist_filtered, imports, dependencies); lower block outputs/import_dependency_unmatched_top5_by_direction.csv and outputs/mentions_alignment_unmatched_top5_by_direction.csv.
% The lower-block mention counts (e.g., code 8,528) come from mentions_alignment_unmatched_top5_by_direction.csv and are correct for "not matched". The 8,530 in mentions_alignment_generic_stoplist_removed_counts.csv counts pairs where the stop list removes the term, a different quantity.
% Stop list: utils.GENERIC_SOFTWARE_NAME_STOPLIST (23 terms).
% top_software_by_view.csv rerun 2026-10-23 under the 23-term list; output identical to the 2026-09-19 file.
\begin{sidewaystable}
\caption{\textbf{Most common software within each view, and software one view carries that another does not.} The upper block gives the most common software names within each view, counted as the number of article-repository pairs in which the name appears, over that view's own eligible population. We show mentions twice, with and without the 23 generic extraction terms in our stop list (code, codes, script, scripts, latex, software, github, gitlab, bitbucket, library, libraries, package, packages, tool, tools, toolbox, pipeline, workflow, api, database, website, server, and notebook). The two mention rows have different eligible populations, because pairs whose only mentions are generic terms drop out once those terms are removed. The lower block gives, for each of the six pairwise directions, the five most common names present in one view and matched by nothing in the other, over pairs published through 2022 for which that comparison is computable; mention names there are shown before the generic-term stop list is applied, to show what mention extraction returns in practice. Names are in their normalized form.}\label{stab:topsoftware}
\footnotesize
\setlength{\tabcolsep}{3pt}
\renewcommand{\arraystretch}{1.1}
\begin{tabular}{@{}>{\raggedright\arraybackslash}p{136pt}>{\raggedleft\arraybackslash}p{36pt}*{5}{>{\raggedright\arraybackslash}p{68pt}}@{}}
\toprule
View or comparison & Eligible pairs & Rank 1 & Rank 2 & Rank 3 & Rank 4 & Rank 5 \\
\midrule
\multicolumn{7}{@{}l}{\textit{Most common within each view}}\\
\addlinespace[2pt]
Mentions (with generic terms) & 55,311 & code (13,717) & scripts (6,624) & github (5,212) & latex (5,040) & matlab (5,014) \\
Mentions (without generic terms) & 52,445 & matlab (5,019) & pytorch (4,467) & linux (2,510) & tensorflow (1,864) & scikitlearn (1,731) \\
Imports & 96,955 & numpy (79,258) & matplotlib (50,308) & scipy (45,278) & torch (41,895) & pandas (35,722) \\
Dependencies & 42,197 & numpy (28,135) & scipy (19,801) & matplotlib (18,549) & pandas (15,530) & tqdm (15,330) \\
\midrule
\multicolumn{7}{@{}l}{\textit{In one view, matched by nothing in the other}}\\
\addlinespace[2pt]
Imports not matched to dependencies & 41,994 & numpy (7,684) & matplotlib (7,018) & scipy (5,804) & torch (5,161) & pandas (4,789) \\
Dependencies not matched to imports & 41,994 & six (7,903) & certifi (7,343) & pythondateutil (7,295) & pyparsing (6,870) & cycler (5,912) \\
\addlinespace
Mentions not matched to imports & 34,014 & code (8,528) & scripts (3,908) & latex (3,687) & github (2,976) & script (2,789) \\
Imports not matched to mentions & 34,014 & numpy (25,187) & matplotlib (15,653) & scipy (14,198) & pandas (12,325) & tqdm (9,448) \\
\addlinespace
Mentions not matched to dependencies & 14,586 & code (3,534) & latex (1,836) & github (1,270) & scripts (1,248) & script (1,004) \\
Dependencies not matched to mentions & 14,586 & numpy (8,971) & scipy (6,331) & matplotlib (5,883) & pandas (5,306) & tqdm (4,963) \\
\botrule
\end{tabular}
\end{sidewaystable}

% Table S6
\begin{table}[p]
\caption{\textbf{Mention rates under both denominators.} Pairs published through 2022. The conditional rate counts only pairs whose article has at least one extracted mention; the unconditional rate counts every pair with the view.}\label{stab:denominators}
\setlength{\tabcolsep}{4pt}
\begin{tabular}{@{}>{\raggedright\arraybackslash}p{51pt}>{\raggedright\arraybackslash}p{85pt}*{5}{>{\raggedleft\arraybackslash}p{34pt}}@{}}
\toprule
View & Denominator & Pairs & Mention rate & Python & R & Pairs naming none \\
\midrule
Imports & Conditional (mention-eligible pairs) & 34,014 & 7.00\% & 5.50\% & 10.81\% & 56.74\% \\
Imports & Unconditional (all pairs with imports) & 97,926 & 2.49\% & 1.83\% & 5.23\% & 84.98\% \\
Dependencies & Conditional & 14,586 & 2.34\% & 2.00\% & 5.18\% & 62.17\% \\
Dependencies & Unconditional (all pairs with dependencies) & 42,615 & 0.77\% & 0.64\% & 2.25\% & 87.05\% \\
\botrule
\end{tabular}
\footnotetext{For the unconditional import rate, the Python and R columns split pairs by the repository's primary language.}
\end{table}

% Table S7
% Source: outputs/mention_predictors_logit_summary.csv (alignment_variant = grouped_hungarian), 16 fits including the four age-only fits.
% Collinearity sentence: outputs/mention_primary_model_collinearity.csv (r = 0.867; VIF 4.22 age, 4.25 popularity, primary model).
\begin{table}[p]
\caption{\textbf{Logistic regression of whether an imported library is mentioned, all sixteen specifications.} We align software names within each pair with the Hungarian algorithm (Methods). Before fitting, we drop libraries whose names match one of 23 generic terms (e.g., \texttt{code}, \texttt{scripts}, \texttt{software}), libraries imported fewer than three times, and pairs above the 99th percentile in number of distinct imported libraries, which leaves 207,026 articles and 22,173 libraries without the cap. Controlled models add article field (the five most common fields plus Other), document type, and publication year as a linear term. Estimates are percentage changes in the odds of a mention per log-unit of popularity or per year of age, with two-way cluster-robust p-values (article and library). ``Cap'' is the 2022 publication-year cap. The popularity coefficient is negative in all twelve specifications that include it (p from $7.3\times10^{-6}$ to 0.198), with p $<$ 0.05 in six, including the primary specification. Age and popularity are strongly correlated (Pearson's r = 0.867), which is consistent with the age estimate turning from negative when fit alone ($-$4.0\% per year with the 2022 cap) to positive once popularity enters the model ($+$7.9\% without controls and $+$3.1\% in the primary specification), although variance inflation factors of 4.2 and 4.3 in the primary specification point to moderate rather than severe collinearity. Full coefficients for the primary specification are in Supplementary Table~\ref{stab:primarymodel}.}\label{stab:logit}
\footnotesize
\setlength{\tabcolsep}{4pt}
\begin{tabular}{@{}llrll@{}}
\toprule
Cap & Model & n (rows) & Popularity: odds change (p) & Age: odds change (p) \\
\midrule
\multicolumn{5}{@{}l}{\textbf{All pairs with imports}}\\
2022 & Age only & 975,239 & -- & $-$4.0\% (0.275) \\
2022 & Popularity only & 975,239 & $-$8.6\% (0.064) & -- \\
2022 & Popularity and age & 975,239 & $-$16.3\% ($<$0.001) & $+$7.9\% (0.015) \\
2022 & Controlled\textsuperscript{\dag} & 975,239 & $-$8.5\% (0.014) & $+$3.1\% (0.505) \\
\addlinespace
None & Age only & 2,493,085 & -- & $-$9.5\% ($<$0.001) \\
None & Popularity only & 2,493,085 & $-$14.8\% ($<$0.001) & -- \\
None & Popularity and age & 2,493,085 & $-$12.9\% (0.007) & $-$1.8\% (0.412) \\
None & Controlled & 2,493,085 & $-$4.8\% (0.198) & $+$0.7\% (0.852) \\
\midrule
\multicolumn{5}{@{}l}{\textbf{Mention-eligible pairs}}\\
2022 & Age only & 346,278 & -- & $-$3.4\% (0.412) \\
2022 & Popularity only & 346,278 & $-$6.8\% (0.197) & -- \\
2022 & Popularity and age & 346,278 & $-$11.7\% (0.003) & $+$4.9\% (0.178) \\
2022 & Controlled & 346,278 & $-$6.7\% (0.069) & $+$2.2\% (0.659) \\
\addlinespace
None & Age only & 385,502 & -- & $-$3.4\% (0.399) \\
None & Popularity only & 385,502 & $-$6.8\% (0.195) & -- \\
None & Popularity and age & 385,502 & $-$11.4\% (0.005) & $+$4.5\% (0.175) \\
None & Controlled & 385,502 & $-$6.5\% (0.077) & $+$1.8\% (0.712) \\
\botrule
\end{tabular}
\footnotetext{\textsuperscript{\dag}Primary specification reported in the main text. --, not in the model. p-values below 0.001 are shown as $<$0.001.}
\end{table}

% Table S8
\begin{table}[p]
\caption{\textbf{ARMM training and evaluation data by seed source and negative strategy.} Positive counts are those of the replication package (164,479 pairs). The 164,572 positives in Methods are the training pipeline's own count, taken before the $\geq$ 2008 publication-year filter (12 pairs) and before one-to-one deduplication against mined pairs added later (81 pairs).}\label{stab:armmdata}
% CHECK: is 160,783 negatives counted before or after deduplication? The train negatives aren't local, and the referee's arithmetic suggests all 6,013 deduplicated rows were negatives.
\begin{tabular}{@{}lrr@{}}
\toprule
Seed source or negative strategy & Pairs & Share \\
\midrule
\textbf{Positive pairs (one-to-one seed pairs), by seed source} & & \\
\quad Papers with Code & 143,434 & 87.2\% \\
\quad SoftCite-2025 & 10,837 & 6.6\% \\
\quad PLOS & 6,886 & 4.2\% \\
\quad JOSS & 2,730 & 1.7\% \\
\quad SoftwareX & 592 & 0.4\% \\
\quad \textbf{Total} & \textbf{164,479} & \textbf{100.0\%} \\
\midrule
\textbf{Negative pairs (synthetic), by strategy} & & \\
\quad Simple shuffle & 16,457 & 10.2\% \\
\quad Random out-of-domain & 16,457 & 10.2\% \\
\quad Nearest-neighbor out-of-domain & 16,453 & 10.2\% \\
\quad Same author, different article & 57,600 & 35.8\% \\
\quad Same contributor, different repository & 53,816 & 33.5\% \\
\quad \textbf{Total} & \textbf{160,783} & \textbf{100.0\%} \\
\botrule
\end{tabular}
\end{table}

% Table S9
% Source: outputs/mention_primary_model_full_coefficients.csv (2026-10-23 refit; popularity and age match the S7 primary row). Reference levels follow statsmodels treatment coding (first level alphabetically): Biochemistry, Genetics and Molecular Biology for field, article for document type.
\begin{table}[p]
\caption{\textbf{Full coefficients for the primary specification.} Logistic regression of whether an imported library is mentioned, fit with controls on all pairs with imports published through 2022 (the primary specification in Supplementary Table~\ref{stab:logit}). Coefficients are on the log-odds scale, with two-way cluster-robust standard errors, 95\% confidence intervals, and p-values (article and library). The model has 975,239 rows from 95,299 articles and 18,281 libraries.}\label{stab:primarymodel}
\footnotesize
\setlength{\tabcolsep}{4pt}
\begin{tabular}{@{}lrrrrr@{}}
\toprule
Term & Coefficient & SE & 95\% CI & p & Odds ratio \\
\midrule
\textbf{Library} & & & & & \\
\quad Popularity (ln cumulative imports) & $-$0.089 & 0.036 & $-$0.161 to $-$0.018 & 0.014 & 0.915 \\
\quad Age (years) & 0.030 & 0.046 & $-$0.059 to 0.120 & 0.505 & 1.031 \\
\addlinespace
\textbf{Field}\textsuperscript{a} & & & & & \\
\quad Computer Science & $-$1.547 & 0.270 & $-$2.077 to $-$1.017 & $<$0.001 & 0.213 \\
\quad Engineering & $-$1.239 & 0.197 & $-$1.624 to $-$0.853 & $<$0.001 & 0.290 \\
\quad Medicine & $-$0.686 & 0.066 & $-$0.816 to $-$0.556 & $<$0.001 & 0.504 \\
\quad Physics and Astronomy & $-$0.592 & 0.113 & $-$0.814 to $-$0.370 & $<$0.001 & 0.553 \\
\quad Other & $-$0.645 & 0.046 & $-$0.736 to $-$0.554 & $<$0.001 & 0.525 \\
\addlinespace
\textbf{Document type}\textsuperscript{b} & & & & & \\
\quad Preprint & $-$0.489 & 0.098 & $-$0.681 to $-$0.298 & $<$0.001 & 0.613 \\
\quad Other & $-$0.333 & 0.073 & $-$0.477 to $-$0.189 & $<$0.001 & 0.717 \\
\addlinespace
\textbf{Article} & & & & & \\
\quad Publication year & 0.005 & 0.024 & $-$0.043 to 0.052 & 0.850 & 1.005 \\
\quad Intercept\textsuperscript{c} & $-$11.277 & 48.288 & $-$105.921 to 83.367 & 0.815 & -- \\
\botrule
\end{tabular}
\footnotetext{\textsuperscript{a}Relative to Biochemistry, Genetics and Molecular Biology. \textsuperscript{b}Relative to journal articles. \textsuperscript{c}Publication year enters in calendar years, so the intercept is the log-odds at year zero and has no direct reading. p-values below 0.001 are shown as $<$0.001.}
\end{table}

\clearpage

%%%%%%%%%%%%%%%%%%%%%%%%%%%%%%%%%%%%%%%%%%%%%%%%%%%%%%%%%%%%%%%%%%%%%%%%%%%%%%%%
% Supplementary Figures
%%%%%%%%%%%%%%%%%%%%%%%%%%%%%%%%%%%%%%%%%%%%%%%%%%%%%%%%%%%%%%%%%%%%%%%%%%%%%%%%

% Fig. S1
\begin{figure}[p]
\centering
\includegraphics[width=\textwidth]{figS1_date_window.pdf}
\caption{\textbf{Days between repository creation and article publication for one-to-one seed pairs.} Publication date minus repository creation date over the 164,479 one-to-one seed pairs in the final data; positive values mean the repository was created before publication. The solid line marks the median (82 days). Dashed lines mark the mining window: repositories created at most 556 days before and 73 days after publication, the 10th and 90th percentiles of one-to-one seed pairs when mining began.}\label{sfig:datewindow}
\end{figure}

% Fig. S2
\begin{figure}[p]
\centering
\includegraphics[width=\textwidth]{figS2_mention_rate.pdf}
\caption{\textbf{Import mention rate by field and article publication year.} Import mention rate among mention-eligible pairs (the conditional rate), by field and article publication year, through 2022; the six most common fields, with field-year cells of fewer than 30 mention-eligible pairs omitted.}\label{sfig:mentionrate}
\end{figure}

% Fig. S3 (print-width redraw deferred)
\begin{figure}[p]
\centering
\includegraphics[width=\textwidth]{figS3_armm_confusion_by_field.pdf}
\caption{\textbf{ARMM confusion matrices on the held-out test set, by field.} Confusion matrices generated from the held-out test set across nine categories: the seven most common named fields, an ``Unknown'' category for articles OpenAlex does not assign a field, and the held-out ``Other'' fields, each under 2\% of the dataset.}\label{sfig:armmconfusion}
\end{figure}

\clearpage

\renewcommand{\refname}{Supplementary References}
\bibliography{main}

\end{document}
